\chapter{结语}


本节不计入考试，但其中一些主题可能对学生仍有参考价值.

\section{现代学科的历史与社会学}

\section{8.1 引言}

在过去三十年左右，代数几何在数学界所处的位置类似于数学在整个世界中的地位，既受尊敬又令人畏惧，远多于被理解.与此同时，我经常被英国同事或华威大学的研究生问到的“服务性”问题通常都很基础:一般来说，这些问题要么在本书中有所涉及，要么在[Atiyah和Macdonald]中有所涵盖.以下内容是对该学科近期发展的一个观点，试图解释这一悖论.我不假装客观.

\section{8.2 史前阶段}

代数几何在19世纪由多个不同来源发展而来.首先是几何传统本身:射影几何(以及当时拿破仑时代军方极感兴趣的描绘几何)、为自身目的研究曲线和曲面、构型几何；然后是复函数理论，将紧致黎曼曲面视为代数曲线，并从其函数域纯代数地重构它.在此基础上，代数曲线与数域整数环之间的深刻类比，以及代数和几何中不变量理论语言的需求，这在20世纪初抽象代数的发展中起到了重要作用.

20世纪的最初几十年见证了深刻的分裂.一方面，是研究曲线和曲面的几何传统，尤其由才华横溢的意大利学派推动；除了自身相当可观的成就外，这一传统在拓扑学和微分几何的发展中起到了重要的激励作用，但越来越依赖于“几何直觉”的论证，甚至连大师们也无法严格维持其严谨性.另一方面，新兴的交换代数力量正在奠定基础并提供证明技术.两种方法差异的一个例子是Chow与van der Waerden之间的争论，后者严格证明了参数化给定次数和属的空间曲线的代数簇的存在，而Severi则终其一生创造性地使用此类参数空间，并在晚年痛恨代数学家(尤其是非意大利人！)闯入他的领域，尤其是暗示他所在学派的工作缺乏严谨性.

\section{8.3 严谨性，第一波浪潮}

继希尔伯特和埃米·诺特引入抽象代数后，van der Waerden、Zariski和Weil在1920年代和1930年代为代数几何奠定了严谨的基础(van der Waerden的贡献常被忽略，显然是因为战后初期包括一些顶尖代数几何学家在内的数学家认为他是纳粹合作者).

他们方案的核心之一是使代数几何能在任意域上工作.在这方面，一个关键的基础性难题是你不能仅仅将簇定义为点集:如果你从一个定义在给定域 \(k\) 上的簇 \(V \subset  {\mathbb{A}}_{k}^{n}\) 开始，那么 \(V\) 不仅仅是 \({k}^{n}\) 的一个子集；你还必须允许 \(V\) 的 \(K\) -值点，针对域扩张 \(k \subset  K\) (参见(8.13, c)的讨论).这也是符号 \({\mathbb{A}}_{k}^{n}\) 的由来，表示簇 \({\mathbb{A}}^{n}\) 的 \(k\) -值点，期望其存在独立于指定域 \(k\) .

允许基域在论证过程中变化的必要性极大地增加了技术和概念上的难度(更不用说符号了).然而，到1950年左右，Weil的基础体系被接受为规范，以至于传统几何学家(如Hodge和Pedoe)感到必须以此为基础编写著作，我认为这大大降低了他们著作的可读性.

\section{8.4 格罗滕迪克时代}

从大约1955年到1970年，代数几何领域由巴黎数学家主导，先是塞尔(Serre)，随后尤其是格罗滕迪克(Grothendieck)及其学派.格罗滕迪克方法的影响力不容低估，尤其是在其方法在某种程度上已不再流行的今天.这一时期取得了巨大的概念和技术进步，得益于方案(scheme，比簇更一般，参见下文(8.12-14))概念的系统应用，代数几何几乎吸收了拓扑学、同调代数、数论等领域的所有进展，甚至在这些领域的发展中发挥了主导作用.格罗滕迪克本人于1970年前后在四十出头时退出了学术舞台，这无疑是一大悲剧性损失(他最初因抗议科学军事资助而离开了IHES).作为一名实践中的代数几何学家，我深刻感受到这一时期发展起来的强大理论工具体系，其中许多仍待以通俗易懂的方式整理成文.

另一方面，格罗滕迪克的人格崇拜带来了严重的副作用:许多毕生致力于掌握魏尔(Weil)基础理论的人遭受了排斥和羞辱，据我所知，只有一两人适应了新的语言；整整一代学生(主要是法国人)被洗脑，愚昧地认为不能用高深抽象形式主义包装的问题不值得研究，因此被排除在数学家自然的发展路径之外——即从一个自己能处理的小问题出发，逐步向外探索.(我实际上知道有一篇关于三次曲面算术的论文最初未被考虑，理由是“构造的自然背景是在一般局部诺特环拓扑空间上”.这绝非玩笑.)当时许多学生似乎没有比研读《EGA》(《代数几何基础》)更高的志向.纯粹为范畴论而学(无疑是所有智力追求中最贫瘠的之一)也始于此时；格罗滕迪克本人不必为此负责，因为他自己对范畴的运用在解决问题上非常成功.

时尚此后又发生了反向转变.在法国最近的一次会议上，我评论了态度的变化，得到的讽刺回答是“但扭曲三次曲线是一个非常好的前代表函子例子”.我了解到，现在负责管理法国科研经费的一些数学家曾在这段知识恐怖时期受苦，因此CNRS科研项目的申请经常被包装得尽量减少与代数几何的关联.

除了极少数能够跟上节奏并存活下来的格罗滕迪克学生外，最持久受益于格罗滕迪克思想并最有效传播这些思想的人，都是间接受其影响的:哈佛学派(通过扎里斯基(Zariski)、芒福德(Mumford)和M.阿廷(M. Artin))，莫斯科的沙法列维奇(Shafarevich)学派，或许还有日本的交换代数学派.

\section{8.5 大爆炸}

历史并未在1970年代初终结，代数几何自那时以来也未曾减少时尚的波动.1970年代，尽管几个大派系各有其特殊兴趣(芒福德与模空间的紧化，格里菲斯(Griffiths)的霍奇理论和代数曲线学派，德利涅(Deligne)与簇的上同调“权重”，沙法列维奇与K3曲面，饭高(Iitaka)及其追随者在高维簇分类等)，我认为我们基本上都相信自己在研究同一学科，代数几何仍是一个整体(实际上还在殖民数学的邻近领域).也许正是因为有一两位专家能够掌握该学科的全部范围，这才成为可能.

到1980年代中期，情况发生了变化，代数几何目前似乎分裂为十多个相互交流有限的学派:曲线与阿贝尔簇，代数曲面与唐纳森理论，三维簇及高维分类，K理论与代数循环，交叉理论与计数几何，一般上同调理论，霍奇理论，特征 \(p\) ，算术代数几何，奇点理论，数学物理的微分方程，弦理论，计算机代数的应用等.

\section{附加脚注与高深评论}

本节混合了基础与高级主题；由于部分内容是对使用本书作为教材的大学教师的“忠告”，或引导高级学生避免学科陷阱，部分材料可能显得晦涩难懂.

\section{8.6 选题}

本书中涉及的主题和示例部分是基于计算难度小和计算简便的实用考虑而选择的.然而，它们也暗示了“簇的分类”:关于圆锥曲线的材料在某种意义上适用于所有有理曲线，而三次曲面是最基本的del Pezzo有理曲面例子.带有群运算的三次曲线是阿贝尔簇(Abelian varieties)的例子；事实(2.2)表明非奇异三次曲线不是有理的，这是分类的第一步.两个平面圆锥曲线的交点(1.12-14)和习题5.6中提到的两个二次曲面的交点也可以纳入类似的模式，其中两个二次曲面的交点在 \({\mathbb{P}}_{k}^{4}\) 中提供了另一类del Pezzo曲面，而两个二次曲面交点上的直线族在 \({\mathbb{P}}_{k}^{5}\) 中提供了二维阿贝尔簇.

曲线的属以及第46页表格中分为三种情况的划分，简而言之就是分类.我本想包含更多关于曲线属的内容，特别是如何用拓扑欧拉示性数或代数几何中的交数来计算属，这些是年轻几何学家必做的五指练习.然而，这些内容足以单独构成一门本科课程，椭圆曲线的复分析理论也是如此.

\section{8.7 计算与理论}

关于这些讲义中的方法，另一点是相当强调那些可以通过显式计算处理的情况.当一般理论证明某种构造的存在时，用显式坐标表达来实现它是一个有益的练习，有助于保持对现实的把握，这也适合本科教材.然而，这不应掩盖理论实际上是为处理复杂情况设计的事实，在那些情况下显式计算往往无法提供任何信息.

\section{8.8 实数域与 \(\mathbb{C}\)}

对实数领域感兴趣的读者可能会失望，因为在 \(§{81} - 2\) 中对 \(\mathbb{R}\) 的处理在 \(§3\) 中让位于对任意域 \(k\) 的讨论，且通常假设该域是代数闭的.我建议这类读者坚持下去；实数几何与复数几何之间存在许多联系，其中一些会令人惊讶.询问实簇的实点是一个非常困难的问题，在代数几何中属于少数兴趣；无论如何，全面了解其复点通常是必要的前提.实数域与 \(\mathbb{C}\) 上的几何之间的另一个直接联系是， \(n\) 维非奇异复簇是 \({2n}\) 维实流形——例如，代数曲面是构造光滑4维流形的主要来源.

除了这些相当明显的联系，还有更微妙的关系，例如:(a) 平面曲线在 \(C \subset  {\mathbb{C}}^{2}\) 上的奇点通过与小球边界的交叉在 \({S}^{3}\) 中产生结；(b) 彭罗斯扭量子构造将带有特殊黎曼度量的4维流形视为一个4维复簇的实值点集，该复簇参数化复3维簇上的有理曲线(因此我们所处的实4维球面 \({S}^{4}\) 可以被识别为复Grassmannian \(\operatorname{Gr}\left( {2,4}\right)\) 中 \({\mathbb{P}}_{\mathbb{C}}^{3}\) 直线的实点集).

\section{8.9 正则函数与层}

已经正确理解了簇 \(X\) 上有理函数 \(f \in  k\left( X\right)\) 的概念以及在点 \(P \in  X\) 处的正则性 \(f\) ((4.7)和(5.4))的读者，对结构层 \({\mathcal{O}}_{X}\) 已有相当直观的认识.对于开集 \(U \subset  X\) ，正则函数的集合是 \(U \rightarrow  k\)

\[
{\mathcal{O}}_{X}\left( U\right)  = \{ f \in  k\left( X\right)  \mid  f\text{ is regular }\forall P \in  U\}  = \mathop{\bigcap }\limits_{{P \in  U}}{\mathcal{O}}_{X,P}
\]

是域 \(k\left( X\right)\) 的子环.层 \({\mathcal{O}}_{X}\) 仅仅是环族 \({\mathcal{O}}_{X}\left( U\right)\) ，当 \(U\) 遍历 \(X\) 的开集时.显然，局部环 \({\mathcal{O}}_{X,P}\) (定义见(4.7)和(5.4))的任一元素在某个包含 \(P\) 的邻域 \(U\) 内是正则的，因此有 \({\mathcal{O}}_{X,P} = \mathop{\bigcup }\limits_{{U \ni  P}}{\mathcal{O}}_{X}\left( U\right)\) .仅此而已；有一个固定的有理截面集合 \(k\left( X\right)\) ，而层在开集 \(U\) 上的截面就是在每个 \(P \in  U\) 处满足正则性条件的有理截面.

该语言足以描述带Zariski拓扑的不可约簇上的任意无挠层.当然，如果 \(X\) 是可约的，或者你想处理更复杂的层，或使用复拓扑，则需要层的完整定义.

\section{8.10 全局定义的正则函数}

如果 \(X\) 是射影簇，那么在每个 \(P \in  X\) 处正则的有理函数 \(f \in  k\left( X\right)\) 只有常数.这是射影簇的一般性质，类似于一复变量函数中的Liouville定理；对于定义在 \(\mathbb{C}\) 上的簇，它源于紧性和最大模原理( \(\left( {X \subset  {\mathbb{P}}_{\mathbb{C}}^{n}}\right.\) 在复拓扑下是紧的，因此全局全纯函数在 \(X\) 上的模必须取得最大值)，但在代数几何中，从头证明这一点出乎意料地困难(参见例如[Hartshorne, I.3.4]；本质上是一个有限性结果，涉及相干上同调群的有限维性).

\section{8.11 射影代数几何的惊人充分性}

Weil对簇的抽象定义(仿射代数集沿同构开集粘合而成)在(0.4)中曾简要提及，并且用层的语言处理起来相当容易.有了这个，单纯考虑嵌入在固定环境空间 \({\mathbb{P}}_{k}^{N}\) 中的簇的想法乍看之下似乎过于限制.我想简要描述一下现代对此问题的看法.

\section{(a) 极化与正性}

首先，簇通常按同构考虑，因此说一个簇 \(X\) 是射影的，意味着 \(X\) 可以嵌入某个 \({\mathbb{P}}^{N}\) ，即同构于如(5.1-7)中所述的闭子簇 \(X \subset  {\mathbb{P}}^{N}\) .拟射影意味着同构于 \({\mathbb{P}}^{N}\) 的局部闭子簇，即射影簇的一个开稠密子集；射影性包含完备性的性质，即 \(X\) 不能作为任何更大簇的稠密开集嵌入.

实际嵌入的选择 \(X \hookrightarrow  {\mathbb{P}}^{N}\) (或非常充足的线丛 \({\mathcal{O}}_{X}\left( 1\right)\) ，其截面将作为 \({\mathbb{P}}^{N}\) 的齐次坐标)通常称为极化，我们用 \(\left( {X,{\mathcal{O}}_{X}\left( 1\right) }\right)\) 表示已做出该选择.除了完备性外，射影簇 \(X \subset  {\mathbb{P}}^{N}\) 满足一个关键的“正次数”条件:若 \(V \subset  X\) 是一个 \(k\) 维子簇，则 \(V\) 与一般线性子空间 \({\mathbb{P}}^{N - k}\) 的交点数为正且有限.反过来，Kleiman判据指出，若一个完备簇上的线丛 \(X\) 在每条曲线 \(C \subset  X\) 上的次数始终大于零(即 \(\geq  \varepsilon\) -(对 \(C\) 的任何合理度量))，则该线丛的某个倍数可用于提供 \(X\) 的射影嵌入.这种正性与复流形上的Kähler度量(与复结构兼容的黎曼度量)选择密切相关.因此，我们将射影性理解为一种“正定性”.

\section{(b) 充分性}

令人惊讶的是，许多代数几何问题都能在射影簇框架内得到解答.前文(8.2)提到的Chow簇构造就是一个例子；另一个是1960年代Mumford的工作，他构造了Picard簇和许多模空间作为拟射影簇(概形).Mori理论(对与有理性相关的簇分类的重要概念进展负责，见[Kollár])是最新的例子；其中的思想和技术本质上不可避免地是射影性的.

\section{(c) 抽象簇的不足}
.
曲线和非奇异曲面自动是拟射影的；但存在非拟射影的抽象簇(奇异曲面，或维数为 \(\geq  3\) 的非奇异簇).然而，如果你需要这些构造，几乎肯定也会需要Moishezon簇(M. Artin的代数空间)，它们是比抽象簇更广义的代数几何对象，通过对“局部片段粘合”的更宽松解释得到.

关于抽象簇的定理通常通过归约到拟射影情形来证明，因此拟射影证明或归约过程的细节哪个更有用、有趣、必要或更可能产生易发表的成果，取决于具体问题和学生的兴趣及就业状况.最近已证明，非奇异的非拟射影抽象簇或Moishezon簇必然包含有理曲线；然而该证明(由J. Kollár完成)属于Mori理论范畴，属于纯正的射影代数几何.

\section{8.12 仿射簇与概形}

代数闭域 \(k\) 上仿射代数簇 \(V\) 的坐标环 \(k\left\lbrack  V\right\rbrack\) (定义4.1)满足两个条件:(i)它是有限生成的 \(k\) -代数；(ii)它是整环.满足这两个条件的环显然是某个簇 \(V\) 的坐标环 \(k\left\lbrack  V\right\rbrack\) ，称为几何环(或几何k-代数).

第二部分有两个关键理论结果；其一是定理4.4，精确说明 \(V \mapsto  k\left\lbrack  V\right\rbrack   = A\) 是仿射代数簇与几何 \(k\) -代数对偶范畴之间的范畴等价(尽管我删去了所有关于范畴的内容，因为不适合年轻读者).另一个是Nullstellensatz定理(3.10)，指出 \(k\left\lbrack  V\right\rbrack\) 的素理想与 \(V\) 的不可约子簇一一对应；而 \(V\) 的点与极大理想一一对应.

综合来看，这些结果将仿射簇(affine varieties) \(V\) 与对应于几何环(geometric rings)的仿射示性层(affine schemes)联系起来(参见定义4.6).

素谱Spec \(A\) 定义于任意环(带1的交换环)为 \(A\) 的素理想集合.它具有Zariski拓扑和结构层；这就是对应于 \(A\) 的仿射示性层(详细内容见[Mumford, Introduction，或Hartshorne，第II章]).仿射示性层比仿射簇更为广泛，有几种截然不同的方式；这些都很重要，我在(8.14)中简要介绍.

理解对于几何环 \(A = k\left\lbrack  V\right\rbrack\) ，素谱Spec \(A\) 包含的信息与簇 \(V\) 完全相同且无多余信息非常重要.NSS告诉我们，点 \(\left. {v \in  V}\right)\) 对应的极大理想 \(\left( {m}_{v}\right.\) 充足，而 \(A\) 的其他素理想 \(P\) 是不可约子簇 \(Y \subset  V\) 上点的极大理想的交:

\[
P = I\left( Y\right)  = \mathop{\bigcap }\limits_{{v \in  Y}}{m}_{v}
\]

忽略簇与示性层的区别是有用且(大致上)允许的，可以用 \(V = \operatorname{Spec}A,v\) 代替 \({m}_{v}\) ，并将素理想 \(P = I\left( Y\right)\) (“一般点”)想象成缝制在子簇 \(Y\) 结构中无处不在的标记.

\section{8.13 意义何在？}

大多数学生除了(8.9)和(8.12)中的内容外，通常不需要了解更多示性层理论，只需注意“泛点”一词在多个技术意义上使用，常常与“足够一般点”含义不同.

本节面向需要研读现代文献的读者，针对示性层理论中各种点的概念提供一些评论，这些概念可能是初学者的主要障碍.

\section{(a) 簇的示性层理论点}

假设 \(k\) 是一个域(可能非代数闭)， \(A = k\left\lbrack  {{X}_{1},\ldots ,{X}_{n}}\right\rbrack  /I\) 与 \(I \subset\)  \(k\left\lbrack  {{X}_{1},\ldots ,{X}_{n}}\right\rbrack\) 为一个理想；记作 \(V = V\left( I\right)  \subset  {K}^{n}\) ，其中 \(k \subset  K\) 为选定的代数闭包.Spec \(A\) 的点比(8.12)中几何环的稍复杂.通过NSS的明显推广， \(A\) 的极大理想由点 \(v = \left( {{a}_{1},\ldots ,{a}_{n}}\right)  \in  V \subset  {K}^{n}\) 确定，即形式为

\[
{m}_{v} = \{ f \in  A \mid  f\left( P\right)  = 0\}  = \left( {{x}_{1} - {a}_{1},\ldots ,{x}_{n} - {a}_{n}}\right)  \cap  A.
\]

很容易看出，不同的点 \(v \in  V \subset  {K}^{n}\) 产生相同的 \(A\) 的极大理想 \({m}_{v}\) 当且仅当它们在伽罗瓦理论意义上是 \(k\) 上的共轭(因为 \(A\) 由系数在 \(k\) 上的多项式组成).因此，极大谱Specm \(A\) 恰好是 \(V\) 的“共轭类”(即 \(\operatorname{Gal}K/k\) 在 \(V\) 上的轨道空间).如(8.12)所示， \(A\) 的其他素理想 \(P\) 对应于一个不可约子簇 \(Y = V\left( P\right)  \subset  V\) (在 \(k\) 上的共轭类)； \(P \in  \operatorname{Spec}A\) 是 \(Y\) 的示性理论上的一般点，也应被视为 \(Y\) 上的一个标记.Spec \(A\) 的Zariski拓扑被调整使得 \(P\) 在 \(Y\) 中处处稠密. \(A\) 的极大理想被称为闭点以示区别.如果 \(C : \left( {f = 0}\right)  \subset  {\mathbb{A}}_{\mathbb{C}}^{2}\) 是一个不可约曲线，它只有一个示性理论上的一般点，对应于 \(\mathbb{C}\left\lbrack  {X,Y}\right\rbrack  /\left( f\right)\) 中的理想(0)，而一个曲面 \(S\) 则在每个不可约曲线 \(C \subset  S\) 中有一个一般点，并且自身还有一个在 \(S\) 中稠密的一般点.

示性理论上的点在将Spec \(A\) 定义为带有拓扑和环层的集合时至关重要(在交换代数中以及代数几何中处理不可约子簇一般点的邻域等概念时也很重要，见(8.14, i))；然而，取值于代数闭包 \(k \subset  K\) 的 \(V \subset  {K}^{n}\) 的点更符合几何意义上的点，称为几何点.这类似于一个簇 \(V\) 的Zariski拓扑更多作为结构层 \({\mathcal{O}}_{V}\) 的载体，而非其自身的几何对象.

\section{(b) 示性理论中的域值点}

若 \(P\) 是 \(A\) 的素理想(即 \(P \in  \operatorname{Spec}A\) 是一个点)，则 \(P\) 处的剩余域是整环 \(A/P\) 的分式域，记为 \(k\left( P\right)\) ；当且仅当 \(P\) 是极大理想时，它是基域 \(k\) 的代数扩张.带有域扩张 \(k \subset  L\) 系数的 \(V\) 的点(点 \(\left( {{a}_{1},\ldots ,{a}_{n}}\right)  \in  V\left( I\right)  \subset  {L}^{n}\) )显然对应于一个同态 \(A \rightarrow  L\) (由 \({X}_{i} \mapsto  {a}_{i}\) 给出)，其核是 \(A\) 的素理想 \(P\) ，或等价地对应于一个嵌入 \(k\left( P\right)  \hookrightarrow  L\) .若 \(P = {m}_{v}\) 是极大理想，且 \(L = K\) 是 \(k\) 的代数闭包，则正是嵌入 \(A/{m}_{v} = k\left( v\right)  \hookrightarrow  K\) 的选择决定了对应于 \(V \subset  {K}^{n}\) 的点的坐标，换言之，区分了该点与其伽罗瓦共轭点.这些即是 \(V\) 的几何点.

对于任意扩张 \(k \subset  L\) ，对应于 \(V\) 的 \(L\) -值点的 \(A \rightarrow  L\) -代数同态可以被修饰得更合理.首先回顾一下，簇不仅仅是点集；即使它只有一个点，也必须说明它定义于哪个域上.因此

\[
\operatorname{Spec}L = \frac{L}{ \cdot  } = {\mathrm{{pt}}}_{L}
\]

是由单个点组成且定义于 \(L\) 上的簇.根据范畴等价(4.4)，态射Spec \(L \rightarrow  V\) (定义于 \(L\) 的点的包含)应当与 \(k\) -代数同态 \(A = k\left\lbrack  V\right\rbrack   \rightarrow  L = k\left\lbrack  {\mathrm{{pt}}}_{L}\right\rbrack\) 相同.

总结方案理论点与域值点之间的关系:点 \(P \in  \operatorname{Spec}A = V\) 是 \(A\) 的一个素理想，因此对应于商同态 \(A \rightarrow  A/P \subset\) Quot \(\left( {A/P}\right)  = k\left( P\right)\) 到一个域.如果 \(L\) 是任意域， \(V\) 的一个 \(L\) -值点是一个同态 \(A \rightarrow  L\) ；方案理论点 \(P\) 以一种套语的方式对应于一个域值点，但域 \(k\left( P\right)\) 随 \(P\) 变化.如果 \(K\) 是 \(k\) 的代数闭包，则 \(V \subset  {K}^{n}\) 的 \(K\) -值点就是几何点；一个 \(K\) -值点 \(v\) 位于一个闭方案理论点 \({m}_{v}\) ，并带有指定的包含 \(A/{m}_{v} = k\left( v\right)  \hookrightarrow  K\) .

\section{(c)魏尔基础中的一般点}

我在(8.3)中提到了魏尔基础中点的特殊性:定义在域 \(k\) 上的一个簇 \(V\) 允许具有任意域扩张 \(k \subset  L\) 上的 \(L\) -值点.这显然源自数论，但在几何中也有影响.例如，如果 \(C\) 是定义在 \(k = \mathbb{Q}\) 上的圆 \({x}^{2} + {y}^{2} = 1\) ，那么

\[
{P}_{\pi } = \left( {\frac{2\pi }{{\pi }^{2} + 1},\frac{{\pi }^{2} - 1}{{\pi }^{2} + 1)}}\right)
\]

被允许作为 \(C\) 的一个 \(\mathbb{C}\) -值点.由于 \(\pi\) 在 \(\mathbb{Q}\) 上是超越的，任何在 \({P}_{\pi }\) 处消失的多项式 \(f \in  \mathbb{Q}\left\lbrack  {x,y}\right\rbrack\) 都是 \({x}^{2} + {y}^{2} - 1\) 的倍数；因此 \({P}_{\pi }\) 是 \(\mathbb{Q}\) -一般点——它不属于任何定义在 \(\mathbb{Q}\) 上的更小子簇.换句话说， \({P}_{\pi }\) 在自同构群Aut \(\mathbb{C}\left( {\left( { \twoheadleftarrow  \operatorname{Gal}\left( {\mathbb{C}/\mathbb{Q}}\right) }\right) \text{ " }}\right)\) 作用下的共轭点在 \(C\) 中稠密.由于 \({P}_{\pi }\) 是 \(\mathbb{Q}\) -一般点，如果你证明了仅涉及 \(\mathbb{Q}\) 上多项式关于 \({P}_{\pi }\) 的命题，那么该命题对 \(C\) 的每个点都成立.

事实上，这一思想已包含在(b)中描述的 \(L\) -值点的概念中，且一般点的几何内涵在此语言中表现得最为清晰.例如，域 \(\mathbb{Q}\left( \pi \right)\) 仅是纯超越扩张，因此 \(\mathbb{Q}\left( \pi \right)  \cong  \mathbb{Q}\left( \lambda \right)\) 和态射Spec \(\mathbb{Q}\left( \lambda \right)  \rightarrow  C\) 是(1.1)中讨论的 \(C\) 的有理参数化:大致上，你可以为超越或未知的 \(\pi\) 代入任何“足够一般”的值.更一般地，有限生成扩张 \(k \subset  L\) 是定义在 \(k\) 上的簇 \(W\) 的函数域；假设 \(\varphi  : \operatorname{Spec}L \rightarrow  V = \operatorname{Spec}A\) 是对应于 \(k\) -代数同态 \(A \rightarrow  L\) 的点，其核为 \(P\) .则 \(\varphi\) 扩展为有理映射 \(f : W \rightarrow  V\) ，其像在子簇 \(Y = V\left( P\right)  \subset  V\) 中稠密，因此 \(\varphi\) 或 \(\varphi \left( {\operatorname{Spec}L}\right)\) 是 \(Y\) 的域值一般点.

\section{(d)方案论中作为态射的点}

(c)中的讨论表明，簇 \(V\) 的一个 \(L\) -值点隐含包含一个有理映射 \(W \rightarrow  V\) ，其中 \(W\) 是与 \(\operatorname{Spec}L\) 双有理的簇(即 \(L = k\left( W\right)\) )；几何学家可以将其视为由 \(W\) 参数化的一族点.

更一般地，对于我们感兴趣的 \(X\) 一个簇(或示性图)，一个 \(S\) 值点(其中 \(S\) 是任意示性图)可以定义为一个态射 \(S \rightarrow  X\) .如果 \(X = V\left( I\right)  \subset  {\mathbb{A}}_{k}^{n}\) 是仿射的，坐标环为 \(k\left\lbrack  X\right\rbrack\) 和 \(S = \operatorname{Spec}A\) ，那么根据(4.4)，一个 \(S\) 值点对应于一个 \(k\) 代数同态 \(k\left\lbrack  X\right\rbrack   \rightarrow  A\) ，即一个满足对所有 \(f \in  I\) 均成立的 \(n\) 元组 \(\left( {{a}_{1},\ldots ,{a}_{n}}\right)\) ，其元素属于 \(A\) ，满足 \(f\left( a\right)  = 0\) .

从高深的角度看，这是簇(variety)概念的最终升华:如果簇 \(X\) 的一个点仅仅是一个态射，那么 \(X\) 本身就是一个函子

\[
S \mapsto  X\left( S\right)  = \{ \text{ morphisms }S \rightarrow  X\}
\]

作用于示性图范畴.(我在第59页脚注中对符号 \({\mathbb{A}}_{k}^{n}\) 的解释已经反映了这一点.)尽管看起来不可思议，这些形而上学的咒语在技术上非常有用，且作为函子定义的簇在现代模空间理论中是基础.给定一个可以“代数地依赖参数”的几何构造(例如固定次数和属的空间曲线)，你可以要求赋予所有可能构造的集合一个代数簇的结构.更进一步，你可以要求存在一个“普遍的”或“包含所有可能构造”的参数空间上的构造族；这个普遍族的参数簇通常可以最直接地定义为一个函子(你仍需证明该簇的存在).例如，(8.2)中提到的Chow簇就代表了该函子

\[
S \mapsto  \{ \text{ families of curves parametrised by }S\} .
\]

\section{8.14 示性图比簇更一般的原因}

我现在单独讨论仿射示性图比仿射簇更一般的三种方式；在严重的情况下，这些复杂情况可能相互叠加，或与(8.11)中讨论的全局问题结合，甚至与诸如 \(p\) -adic收敛或Arakelov赫尔米特度量等新现象结合.空间的考虑幸运地使我无需对这些迷人话题做更多阐述.

\section{(i) 不局限于有限生成代数}

假设 \(C \subset  S\) 是一个非奇异仿射曲面上的曲线(如果必须，基域为 \(\mathbb{C}\) ).环

\[
{\mathcal{O}}_{S,C} = \{ f \in  k\left( S\right)  \mid  f = g/h\text{ with }h \notin  {IC}\}  \subset  k\left( S\right)
\]

是 \(S\) 在 \(C\) 处的局部环；元素 \(f \in  {\mathcal{O}}_{S,C}\) 在包含 \(C\) 的一个稠密开子集的开集上是正则的.该环中的可除性理论非常优美，并且与亚纯函数的零点和极点的几何概念相关: \(C\) 局部由单个方程 \(\left( {y = 0}\right)\) 定义， \(y \in  {I}_{C}\) 是局部生成元，每个非零元素 \(f \in  {\mathcal{O}}_{S},C\) 都可写成 \(f = {y}^{n} \cdot  {f}_{0}\) 的形式，其中 \(n \in  \mathbb{Z}\) ，且 \({f}_{0}\) 是 \({\mathcal{O}}_{S,C}\) 的可逆元.具有此性质的环称为离散赋值环(discrete valuation ring，d.v.r.)，以离散赋值 \(f \mapsto  n\) 命名，该赋值计数 \(f\) 沿 \(C\) 的零点阶数( \(n < 0\) 对应极点)；元素 \(y\) 称为 \({\mathcal{O}}_{S,C}\) 的局部参数.

现在，层理论允许我们大胆地将 Spec \({\mathcal{O}}_{S,C}\) 视为一个几何对象，即只有两个点的拓扑空间:一个闭点，极大理想 \(\left( y\right) \left( { = \text{ the generic point of }C}\right)\) ，和一个非闭点，零理想 \(0\left( { = \text{ the generic point of }S}\right)\) .这里的优势不仅仅是技术性的:离散赋值环的简单交换代数当然早已被用来证明代数几何和复函数理论中的结果(例如，关于函数的理想，或关于分支覆盖 \(T \rightarrow  S\) 在 \(C\) 上方的局部行为，借助域扩张 \(k\left( S\right)  \subset  k\left( T\right)\) 来描述)——这些都早于层的发明.更重要的是，它为我们提供了精确的几何语言和局部代数的简明图景.

以上只是一个例子，涉及局部化，或“子簇的典型点邻域”的概念，说明了从取比域上有限生成代数更一般的环的 Spec 对普通几何的益处；类似的例子是将簇 \(W\) 的典型点 Spec \(k\left( W\right)\) 看作是 \(W\) 所有非空开集的交集所得到的簇(参见(8.13, c))，就像柴郡猫的脸消失后留下的笑容.

\section{(ii) 拟零元}

环 \(A\) 可以包含拟零元；例如 \(A = k\left\lbrack  {x,y}\right\rbrack  /\left( {{y}^{2} = 0}\right)\) 对应于“双重直线” \(2\ell  \subset  {\mathbb{A}}_{k}^{2}\) ，可被视为直线的无穷小条带邻域. \(A\) 中的元素形如 \(f\left( x\right)  + \varepsilon {f}_{1}\left( x\right)\) (其中 \({\varepsilon }^{2} = 0\) )，因此它看起来像是在 \(\ell\) 处截断到一阶的多项式泰勒级数展开.如果你每天多次刻意练习，应该能够将其形象化为双重直线 \(2\ell\) 上的函数.

拟零元使层理论能够处理任意阶截断的泰勒级数，例如通过幂级数方法处理簇上的点.它们在(8.14, d)末尾讨论的模问题中至关重要:例如，它们为处理几何构造的一阶无穷小变形(作为参数空间 \(\operatorname{Spec}\left( {k\left\lbrack  \varepsilon \right\rbrack  /\left( {{\varepsilon }^{2} = 0}\right) }\right)\) 上的构造)提供了精确语言，并将这些视为普适参数簇的切向量.它们还开启了一系列无经典对应的现象，例如特征 \(p\) 下不可分域扩张与簇上向量场的李代数之间的关系.

\section{(iii) 无基域}

设 \(p\) 为素数， \({\mathbb{Z}}_{\left( p\right) } \subset  \mathbb{Q}\) 为分母中不含 \(p\) 的有理数子环； \({\mathbb{Z}}_{\left( p\right) }\) 是另一个带参数 \(p\) 的离散赋值环.它有唯一的极大理想 \(0 \neq  p{\mathbb{Z}}_{\left( p\right) }\) ，其剩余域为 \({\mathbb{Z}}_{\left( p\right) }/p{\mathbb{Z}}_{\left( p\right) } \cong  {\mathbb{F}}_{p} = {\mathbb{Z}}_{\left( p\right) }/\left( p\right)\) .若 \(F \in  {\mathbb{Z}}_{\left( p\right) }\left\lbrack  {X,Y}\right\rbrack\) ，则考虑曲线 \({C}_{\mathbb{C}} : \left( {F = 0}\right)  \subset  {\mathbb{A}}_{\mathbb{C}}^{2}\) 是有意义的，或者取 \(F{\;\operatorname{mod}\;p}\) 的约化 \(f\) ，并考虑曲线 \({C}_{p} : \left( {f = 0}\right)  \subset  {\mathbb{A}}^{2}{\mathbb{F}}_{p}\) .包含复数域上的曲线和有限域上的曲线的几何对象是什么？是否将其视为真正的几何对象见仁见智，但概形 \(\operatorname{Spec}{\mathbb{Z}}_{\left( p\right) }\left\lbrack  {X,Y}\right\rbrack  /\left( F\right)\) 正是这样做的.

这同样从技术上讲并非新思想:自18世纪以来，模 \(p\) 约化曲线的做法已被采用，魏尔(Weil)基础理论中包含了处理此问题的“专门化”理论.其优势在于对离散赋值环(d.v.r.) \({\mathbb{Z}}_{\left( p\right) }\) 上的曲线 \(\operatorname{Spec}{\mathbb{Z}}_{\left( p\right) }\left\lbrack  {X,Y}\right\rbrack  /\left( F\right)\) 作为纤维于 \(\operatorname{Spec}\left( {\mathbb{Z}}_{\left( p\right) }\right) \left( {{}^{\iota } = \left( {\cdot  - }\right) {}^{\iota }}\right)\) 上的几何对象有更好的概念性理解，其中两条曲线 \({C}_{\mathbb{C}}\) 和 \({C}_{p}\) 分别作为一般纤维和特殊纤维.

同理，对于 \(F \in  \mathbb{Z}\left\lbrack  {X,Y}\right\rbrack\) ，概形 \(\operatorname{Spec}\mathbb{Z}\left\lbrack  {X,Y}\right\rbrack  /\left( F\right)\) 是一个几何对象，包含对每个素数 \(p\) 定义在 \({\mathbb{F}}_{p}\) 上的曲线 \({C}_{p} : \left( {{f}_{p} = 0}\right)  \subset  {\mathbb{A}}_{{\mathbb{F}}_{p}}^{2}\) ，其中 \({f}_{p}\) 是 \(F{\;\operatorname{mod}\;p}\) 的约化，同时也包含曲线 \({C}_{\mathbb{C}} : \left( {F = 0}\right)  \subset  {\mathbb{A}}_{\mathbb{C}}^{2}\) ，称为算术曲面；它还包含许多其他内容:特别地，对于每个坐标为代数数的点 \(c \in  {C}_{\mathbb{C}}\) ，它包含一份 \(\operatorname{Spec}\mathbb{Q}\left\lbrack  c\right\rbrack\) ，因此基本上包含了定义曲线 \(c\) 的数域 \(\mathbb{Q}\left( c\right)\) 的整数环的全部信息.

尽管这个对象乍一看极其荒诞不经(如果你多加练习，还是能习惯的)，它却是现代数论的关键组成部分，也是Arakelov(阿拉凯洛夫)和Faltings(法尔廷斯)工作的基本基础.

\section{8.15 三次曲面上直线存在性的证明}

每个成熟的代数几何学家都知道通过维数计数证明(7.2)的传统方法(例如参见[Beauville，《复代数曲面》，第50页]，或[Mumford，《代数几何I，复射影簇》，第174页]).我先简要回顾这一证明，然后再讨论其中的难点.

\({\mathbb{P}}^{3}\) 上的直线集合由4维Grassmann流形 \(\operatorname{Gr} = \operatorname{Gr}\left( {2,4}\right)\) 参数化，三次曲面由三次型在(X, Y, Z, T)上的射影空间 \(S = {\mathbb{P}}^{N}\) 参数化(实际上是 \(N = {19}\) ).记 \(Z \subset  \operatorname{Gr} \times  S\) 为对应的关联子流形.

\[
Z = \{ \left( {\ell ,X}\right)  \mid  \ell  \in  \operatorname{Gr},X \in  S\text{ s.t. }\ell  \subset  X\} .
\]

由于在给定直线 \(\ell\) 上消失的三次型构成一个 \({\mathbb{P}}^{N - 4}\) ，从第一射影 \(Z \rightarrow\) Gr可以容易推断出 \(Z\) 是一个有理的 \(N\) 维流形.因此第二射影 \(p : Z \rightarrow  S\) 是两个 \(N\) 维流形之间的态射，因此

(i) 要么像 \(p\left( Z\right)\) 的像是 \(S\) 中一个 \(N\) 维流形，因此包含 \(S\) 的一个稠密开集，要么 \(p\) 的每个纤维维数为 \(\geq  1\) .

(ii) \(Z\) 是射影流形，因此像 \(p\left( Z\right)\) 在 \(S\) 中是闭集.

由于存在仅包含有限条直线的三次曲面，(i)中的第二种可能性不成立，因此每个足够一般的三次曲面都包含直线.然后(ii)保证了 \(p\left( Z\right)  = S\) ，从而每个三次曲面都包含直线.

我认为这个论证不适合本科课程，原因有二:(i)的陈述假设了关于纤维维数的结果，虽然直观上可接受(尤其是课程最后一周的学生)，但严格证明较难；而(ii)是射影流形完备性的定理，这同样需要证明(通过消元理论、紧性，或完备性值准则的完整论述).

据我所知，我在(7.2)中的证明是新的；有见识的读者当然会看到它与通过向量丛的传统论证的联系:Grassmann流形 \(\operatorname{Gr}\left( {2,4}\right)\) 具有一个典范的秩为2的向量丛 \(E\) (由 \({\mathbb{P}}^{3}\) 上直线的线性形式组成)；将三次曲面方程 \(f\) 限制到每条直线 \(\ell  \subset  {\mathbb{P}}^{3}\) 上定义了 \(E\) 的三次对称幂的截面 \(s\left( f\right)  \in  {S}^{3}E\) .最后，每个 \({S}^{3}E\) 的截面必有零点，这要么由 \(E\) 的放大性，要么由Chern类论证保证(后者也给出了神奇的数字27).

\section{代序言}

\section{8.16 致谢与名人提及}

试图列出所有对我教育有贡献的数学家是徒劳的.我深深感激我的正式导师Pierre Deligne(皮埃尔·德林格)和Peter Swinnerton-Dyer(彼得·斯温纳顿-戴尔)(在他成为成功的政治家和媒体人物之前)；我可能从David Mumford(大卫·芒福德)的著作中学到最多，我对Grothendieck(格罗滕迪克)遗产的理解(无论如何)主要来自Mumford和Deligne.作为数学家和作为一个人，我的世界观深受Andrei Tyurin(安德烈·图林)的影响.

我对本科代数几何课程应有内容的看法部分基于彼得·斯温纳顿-戴尔(Peter Swinnerton-Dyer)于1970年左右为剑桥三一学院设计的一门课程，随后由他和巴里·特尼森(Barry Tennison)教授；我的书在某些方面是这门课程的直接传承，其中一些习题完全照搬了特尼森的习题册.然而，我从华威大学课程结构所允许的自由中获益匪浅，特别是教学理念(由克里斯托弗·齐曼(Christopher Zeeman)明确提出)——研究经验必须作为决定如何以及教授什么内容的主要指导原则.

出现在(2.14-)中的“眨眼环面”来自吉姆·伊尔斯(Jim Eells)，他告诉我这是从H. 霍普夫(H. Hopf)那里学来的(这可能源自德国较早的数学艺术传统).我必须感谢卡罗琳·西里斯(Caroline Series)、弗兰斯·奥尔特(Frans Oort)、保罗·科恩(Paul Cohn)、约翰·琼斯(John Jones)、乌尔夫·佩尔森(Ulf Persson)、大卫·福勒(David Fowler)、一位匿名审稿人以及剑桥大学出版社(C.U.P.)的大卫·特拉纳(David Tranah)对本书预印本提出的宝贵意见，并为在某些情况下未能完全采纳他们的建议或坚持己见表示歉意.

我感谢马蒂娜·耶格尔(Martina Jaeger)对第一版的多处修正，以及若林功(Isao Wakabayashi)详尽的审读，发现了许多不准确之处.特别感谢R.J. 查普曼(R.J. Chapman)和比尔·布鲁斯(Bill Bruce)指出第一版中最严重的错误(我在(7.2)开头避免提及黑塞矩阵(Hessian)，是通过引用一个错误的命题作为习题来掩盖的).

\section{旧索引}

抽象簇 \(4,{79} - {80},{119} - {120}\)

升链条件(ascending chain condition，a.c.c.) \({48} - {49},{53},{55},{63}\)

仿射坐标变换 12,24

射影簇上的仿射锥 81,82

仿射坐标 14,38,43,83,112

射影簇的仿射覆盖 83

仿射曲线 39,45,79

射影簇的仿射片 \({13},{38},{79},{83} - {84},{92},{111}\)

仿射示性空间 120

仿射空间 \({\mathbb{A}}_{k}^{n}{50},{53},{60},{64},{66},{69},{77},{79},{89},{94},{100},{115},{123}\)

仿射簇 \(4,{50},{70} - {71},{72},{74},{78},{120}\)

代数(子)集 \({50} - {55},{64},{66},{78},{81} - {84}\)

代数闭域 52, 54, 55, 64, 71, 77, 118, 120, 121, 122

代数无关 59, 89, 97, 101

渐近线 9, 12, 14, 60, 112

B Bézout定理 \({17} - {18},{33},{35} - {36},{112}\)

双有理等价 \({87} - {89},{99},{100} - {101},{107} - {108},{123}\)

双有理映射 \({87} - {89},{91},{92},{99},{100} - {101}\)

爆破 100-101

C 几何范畴 \(2 - 4,{46}\)

范畴论 \(4,{116},{120},{123} - {124}\)

特征 \({p4},{14},{16},{24},{28},{61} - {62},{107},{125}\)

簇的分类 \({43} - {47},{117}\)

完备簇 119

复分析几何 \(3,{36},{43} - {47},{95},{118},{119}\)

复变函数论 \(6,{45} - {47},{114},{118}\)

圆锥曲线 \(9 - {21},{25},{30} - {33},{37} - {38},{45},{85},{93},{106}\)

坐标环 \(k\left\lbrack  V\right\rbrack  {66} - {72},{73},{74},{75},{120},{123}\)

\(V - I\) 对应关系 \({50} - {51},{52},{53},{54},{55},{60},{63} - {64},{66} - {67},{81} - {82},{84},{120},{121}\)

三次曲线 \(1,2,7,{27} - {42},{43} - {44},{75} - {77},{79},{92},{117}\)

三次曲面 \(6,{102} - {113},{116},{117},{126}\)

尖点三次曲线27,41,68,74,103,112

有理函数的分母D \(4,{68},{72},{76} - {77},{78}\)

稠密开集36,51,67,71,72,73,88,95,97,99

维数2,57,59,60,62,64,97,99,102,118,126

丢番图问题 \(1,9 - {10},{24},{28},{41} - {42},{45} - {47},{125} - {126}\)

离散赋值环(discrete valuation ring, d.v.r.)124, 125

判别式22, 23, 106-107

定义域 dom \({f71} - {73},{77},{78},{83} - {84},{85},{87},{91}\)

占优映射 \({73} - {74},{87}\)

消元理论 \({25} - {26},{57},{64},{104},{105} - {106},{113}\)

空集 \(\varnothing 0,{45},{52},{53},{55},{73},{82}\)

范畴等价 \(V \mapsto  k\left\lbrack  V\right\rbrack  {69},{120},{122}\)

欧拉公式100,111

有限代数 4, 57-58, 59, 60, 61, 64,

有限生成代数 \(4,{49},{54},{57} - {59},{71},{120},{124}\)

有限生成理想48,49,50,81

型 \({16} - {17},{22},{25},{30},{99}\)

函数域 \(k\left( V\right) {62} - {63},{71},{73},{74},{7883},{85},{87},{88},{89},{97},{99},{114},{123},{124}\)

典型点 121, 122-123, 124

曲线的属 \({43} - {47},{115}\)

三次曲线上的群运算 \({33} - {36},{39} - {41},{46},{76}\)

黑塞矩阵(Hessian), 39, 103, 111-112, 127

齐次理想 \({80} - {81},{84}\)

齐次多项式 \(\left( { = \text{ form }}\right) {16} - {17},{22},{25},{30},{80} - {81},{99} - {100}\)

齐次 \(V - I\) 对应 \({81} - {82}\)

超曲面 50, 56-57, 62-63, 64, 88-89, 94-95, 99, 101

无限下降法(infinite descent) 29, 42

拐点 34, 38-39, 41, 103, 112

平面曲线的交点 17, 33, 35-36, 64

两圆锥曲线的交点 \({20} - {25},{117}\)

两个二次曲面的交点 \({91} - {92},{117}\)

不可约代数集 \({33},{52} - {53},{55},{57},{63},{67},{71},{78},{82},{84},{92},{95}\)

不可约超曲面 \({56} - {57},{64}\)

同构 4, 68, 70, 74-75, 77, 78, 79, 85, 87, 90, 92, 93, 99

雅各布森环(Jacobson ring) 121

笑话(非考试内容) 51, 55, 69, 91, 116

平面曲线的线性系统 \({18} - {20},{30} - {33}\)

线性射影 \({10},{60},{65},{68},{86},{92},{107} - {108}\)

局部环 \({\mathcal{O}}_{V,P}{71},{83},{118},{124}\)

局部化 \(A\left\lbrack  {S -  - 1}\right\rbrack  {49},{56},{63},{72},{124}\)

极大理想 \({m}_{P}{54},{55},{64},{120},{121},{122}\)

军事资金 13, 114, 115

模空间(moduli)46,47,116,123,125

单项式曲线 27, 57, 64

态射(= 正则映射) 4, 36, 74, 76, 77, 80, 85, 90, 93, 108, 112, 123 重根，重数 16-17, 34, 35, 38, 40, 52, 94, 102, 107

N 个结点的三次曲线 27, 40, 68, 78, 103, 113

诺特正规化(Noether normalisation) \({59} - {63},{64}\)

Zariski 拓扑的诺特性质(Noetherian property) 53

诺特环(Noetherian ring) \({48} - {49},{63}\)

非奇异 2, 33, 92, 94-95, 97, 99, 101, 102, 107, 111, 112, 113, 118, 124

非奇异三次曲线，见三次曲线

三次曲线的标准形 \({38} - {40},{41}\)

零点定理(Nullstellensatz) \(4,{30},{54},{72},{81} - {82},{120} - {121}\)

数论，见丢番图问题

O 开集，见稠密开集或标准开集

P 平行性 11, 12, 14, 15, 25, 60

参数化曲线 9-10, 15-16, 17-18, 24, 27-28, 31, 40, 45, 47, 68, 74, 77-78, 85, 86, 88, 123

帕斯卡神秘六边形 36-37

圆锥曲线束 \({20},{21} - {25}\)

无穷远点9,12,13,14,16,17,38,39,40,43,60,76,112

极线 104, 113

多项式函数 \(2,3,4,{51},{66} - {70},{72},{96}\)

多项式映射 2, 67-70, 74, 77, 78

素理想52,55,60,61,120,122

素谱 Spec \({A120},{121},{124}\)

本原元定理 62

主理想整环 (PID) 63

簇的乘积78,89,92

射影代数几何 119-120

射影坐标变换 13,41

射影曲线13,24,44,75

射影等价13,15,18

射影几何9,11,79

射影直线 \({\mathbb{P}}^{1}{16},{43},{79},{80},{85},{90}\)

射影平面 \({\mathbb{P}}^{2}9,{11} - {20},{17},{25},{30} - {33},{38},{47},{79},{86},{107}\)

射影空间 \({\mathbb{P}}^{3},{\mathbb{P}}^{n}4,6,{60},{80},{81},{85},{86},{89},{91},{99} - {100},{102},{108}\)

射影簇 \(4,{13},{79} - {90},{119}\)

射影簇与非奇异性 \({99} - {100}\)

Q 二次曲面 64, 86, 90-91, 91, 108, 109, 111, 113

拟射影簇 4, 119

\(\mathbf{R}\) 根(radical) rad \(I = \sqrt{I}{54} - {55},{63},{81} - {82}\)

有理曲线 45,85,91,117,120

有理函数 \(3,4,{28},{45},{68},{71} - {72},{76},{78},{82} - {83},{118}\)

有理映射 4, 28, 72-74, 76-78, 84-88, 91, 107-108, 112, 123

有理正规曲线 85, 91

有理簇 45,88,107,115,126 实几何 \(6 - 7,{117}\)

正则函数 \(2,4,{71} - {72},{77},{78},{118},{124}\)

射影簇上的正则函数 \({80},{82},{83},{90},{92},{118} - {119}\)

正则映射(=态射) 2,4,6,71-72,74,77,78,85,90,92,112

结果式 25, 64, 103-106, 113

黎曼球面 43

黎曼曲面 43,45,112

二元形式的根，见零点

S Segre 嵌入 89

可分性 \({61} - {62},{95},{125}\)

奇异 2,94,95,97,100,101,102,111,112

奇异圆锥曲线 \({21} - {22},{25},{106},{112}\)

奇异三次曲线，见结点三次曲线或尖点三次曲线

奇异三次曲面 112

奇点 \(2,7,{28},{91},{94},{100} - {101},{111},{118}\)

奇点理论 \(2,6,{100} - {101},{117}\)

标准仿射片 \({V}_{\left( i\right) }{13},{38},{79},{83} - {84},{92}\)

标准开集 \({V}_{f}{55},{72},{74} - {75},{98}\)

T 切空间 \({T}_{P}{V2},{33},{34},{40},{41},{94} - {101},{102},{111},{125}\)

曲线的拓扑 \({43} - {45},{46}\)

拓扑，参见Zariski拓扑

超越次数 tr \({\deg }_{k}{K63},{88},{89},{97},{101}\)

直线的横截线108,110,113

扭曲三次曲线85,91,116

U 唯一分解整环 (UFD) 28, 54, 63, 71, 78

V 丛 \({50},{57},{70} - {71},{80},{88},{89},{97},{99},{102},{115},{118},{119} - {124}\)

Veronese曲面 93

Z Zariski拓扑 36, 50-51, 64, 67, 71, 73, 75, 78, 81, 83, 84, 89, 92, 95, 118, 120, 121-122

型的零点 \({16} - {17},{22},{23},{25},{31},{34},{38},{41},{103},{107},{113}\)
